% =============================================================
% CRS Ablation Study — LaTeX Tables & Method Description
% Experiment: hf_crs_ablation_rtx_pro_10939241
% Requires: \usepackage{booktabs, multirow, graphicx, pgfplots}
% =============================================================

% ----- Method Description -----
\subsection{Cross-Round Sharing (CRS)}
\label{sec:crs}

In the standard KV-reuse pipeline, each agent must perform at least one
dense prefill for every new user question before its KV cache can be
reused by later requests.  Specifically, when an anchor entry exists for
a given placeholder but the current agent has not yet stored its own
context delta, the system forces a dense prefill to compute that delta.
This creates a cold-start bottleneck: for the first occurrence of each
input message, \emph{every} agent falls back to dense prefill regardless
of available anchor history.

Cross-Round Sharing (CRS) eliminates this bottleneck by allowing
downstream agents (e.g., Agent~2, Agent~3, and the FinalRefer aggregator)
to borrow the context delta computed by an upstream agent (Agent~1) that
has already completed dense prefill for the same message within the
current round.  When the \texttt{has\_active\_anchor} gate detects that
(i)~an anchor entry exists, (ii)~the current agent lacks its own delta,
and (iii)~an upstream agent's delta is available, CRS \emph{bypasses} the
forced dense prefill and instead proceeds with KV reuse.  During the
reuse phase, CRS applies a weighted cross-delta correction: using the
upstream agent's stored key/value deltas and historical offset statistics,
it estimates the offset that would have been obtained by a full dense
prefill.  The correction weights are derived from the L2-norm similarity
between the upstream delta and historical delta pairs across a sliding
window.

\subsubsection{Formulation}

Each agent stores its anchor entries as an absolute delta from the
rotated base cache.  For a shared token sequence $x$ at placeholder $p$,
when Agent~1 (upstream) has already completed dense prefill this round:
%
\begin{equation}
  \text{Agent 1 stores:}\quad \Delta_{\text{base}\to 1} = \mathrm{KV}(x \mid P_1) - \mathrm{KV}(x \mid \text{base})
  \label{eq:crs-delta1}
\end{equation}
%
\begin{equation}
  \text{Agent 2 stores:}\quad \Delta_{\text{base}\to 2} = \mathrm{KV}(x \mid P_2) - \mathrm{KV}(x \mid \text{base})
  \label{eq:crs-delta2}
\end{equation}
%
Both deltas are already present in the anchor pool from prior requests.
Agent~2 needs $\mathrm{KV}(x \mid P_2)$ for the \emph{current} input
but has not yet run dense prefill, so its own
$\Delta_{\text{base}\to 2}$ is unavailable this round.  CRS
approximates it by rewriting:
%
\begin{equation}
  \mathrm{KV}(x \mid P_2) = \mathrm{KV}(x \mid P_1) + \underbrace{\Delta_{\text{base}\to 2} - \Delta_{\text{base}\to 1}}_{\displaystyle \Delta_{1\to 2}}
  \label{eq:crs-rewrite}
\end{equation}
%
Since Agent~1's current-round cache $\mathrm{KV}(x \mid P_1)$ is
available, the problem reduces to estimating the cross-agent correction
$\Delta_{1\to 2}$.  We approximate it as a weighted average over $W$
historical anchor entries where \emph{both} agents have stored deltas:
%
\begin{equation}
  \Delta_{1 \to 2} \approx \sum_{i=1}^{W} w_i \bigl(\Delta_{\text{base}\to 2}^{(i)} - \Delta_{\text{base}\to 1}^{(i)}\bigr)
  \label{eq:crs-cross}
\end{equation}
%
The weights $w_i$ are computed from the L2 similarity between the
current base key cache and each historical anchor's key embedding:
%
\begin{equation}
  w_i = \frac{\exp\!\bigl(-\|K^{\text{base}} - K^{(i)}_{\text{emb}}\|_2\bigr)}{\sum_{j=1}^{W} \exp\!\bigl(-\|K^{\text{base}} - K^{(j)}_{\text{emb}}\|_2\bigr)}
  \label{eq:crs-weights}
\end{equation}
%
Substituting back, the final CRS estimate becomes:
%
\begin{equation}
  \mathrm{KV}(x \mid P_2) \approx \mathrm{KV}(x \mid \text{base}) + \Delta_{\text{base}\to 1} - \sum_{i=1}^{W} w_i \bigl(\Delta_{\text{base}\to 1}^{(i)} - \Delta_{\text{base}\to 2}^{(i)}\bigr)
  \label{eq:crs-final}
\end{equation}
%
When no historical pairs exist ($W = 0$), the correction vanishes and
CRS reduces to directly reusing Agent~1's delta.  The same procedure is
applied independently to value caches.

% ----- Experimental Setup Table -----
\begin{table}[t]
\centering
\caption{Experimental setup for the CRS ablation study.}
\label{tab:crs-setup}
\renewcommand{\arraystretch}{1.15}
\begin{tabular}{|l|l|}
\hline
\textbf{Parameter}        & \textbf{Value} \\
\hline\hline
Model                     & Llama-3.1-8B-Instruct \\
\hline
Benchmark                 & GSM8K (1{,}319 questions) \\
\hline
GPU                       & NVIDIA RTX PRO 6000 \\
\hline
Backend                   & HF (\texttt{gpt\_chat}) \\
\hline
Topology                  & FullConnected, 3$\times$MathSolver + FinalRefer \\
\hline
KV threshold              & 0.3 \\
\hline
Max anchors               & 20 \\
\hline
Window size               & 5 \\
\hline
Temperature               & 1.0 \\
\hline
Execution mode            & \texttt{allow\_kv\_reuse} \\
\hline
\end{tabular}
\end{table}

% ----- Main Results Table -----
\begin{table}[t]
\centering
\caption{CRS ablation: KV reuse with and without Cross-Round Sharing.
Both runs use the same code version and hyperparameters on the full
GSM8K test set (1{,}319 questions).  The baseline disables CRS via
\texttt{-{}-no-current-round-sharing}.}
\label{tab:crs-results}
\renewcommand{\arraystretch}{1.15}
\begin{tabular}{|l|c|c|}
\hline
                          & \textbf{CRS OFF (Baseline)} & \textbf{CRS ON} \\
\hline\hline
Final Accuracy            & 0.8218                      & 0.7726 \\
\hline
KV Reuse Rate             & 70.2\%                      & 83.6\% \\
\hline
KV Reuse Calls            & 3{,}703 / 5{,}276           & 4{,}410 / 5{,}276 \\
\hline
CRS Bypass (dense$\to$reuse) & ---                      & 624 \\
\hline
Reuse Rate w/o CRS Bypass & ---                         & 71.8\% \\
\hline
\end{tabular}
\end{table}

% ----- Per-Node Accuracy Table -----
\begin{table}[t]
\centering
\caption{Per-node accuracy and mode distribution for CRS OFF (baseline) vs.\ CRS ON.
  Node~1 is the upstream agent and does not receive CRS corrections.}
\label{tab:crs-per-node}
\renewcommand{\arraystretch}{1.15}
\begin{tabular}{|l|cc|cc|cc|}
\hline
                     & \multicolumn{2}{c|}{\textbf{Accuracy}} & \multicolumn{2}{c|}{\textbf{Dense Prefill}} & \multicolumn{2}{c}{\textbf{KV Reuse}} \\
\cline{2-7}
\textbf{Node (Role)} & OFF   & ON    & OFF & ON  & OFF & ON \\
\hline\hline
Node 1 (Math Solver)    & 77.26\% & 77.48\% & 280 & 258 & 1{,}039 & 1{,}061 \\
\hline
Node 2 (Math Analyst)   & 37.30\% & 39.04\% & 385 & 119 &   934 & 1{,}200 \\
\hline
Node 3 (Programmer)     & 74.45\% & 65.88\% & 405 & 161 &   914 & 1{,}158 \\
\hline
Node 0 (FinalRefer)     & 82.18\% & 77.26\% & 503 & 328 &   816 &   991 \\
\hline
\end{tabular}
\end{table}

% ----- CRS Applied vs Not — Accuracy Table -----
\begin{table}[t]
\centering
\caption{Final accuracy (FinalRefer, model-evaluated) on questions where
  CRS was applied vs.\ not (CRS~ON run, 1{,}319 questions).  CRS was
  triggered on 228 questions that had upstream delta available.}
\label{tab:crs-applied-vs-not}
\renewcommand{\arraystretch}{1.15}
\begin{tabular}{|l|c|c|c|}
\hline
                           & \textbf{CRS Applied}        & \textbf{No CRS}             & \textbf{$\Delta$} \\
                           & (228 questions)              & (1{,}091 questions)         &                   \\
\hline\hline
FinalRefer Accuracy        & 80.70\%                      & 76.54\%                     & $+4.17$           \\
\hline
\end{tabular}
\end{table}

% ----- CRS Distribution by Node -----
\begin{table}[t]
\centering
\caption{Distribution of CRS corrections across downstream nodes.
  Node~1 serves as the upstream delta source and is never corrected.}
\label{tab:crs-distribution}
\renewcommand{\arraystretch}{1.15}
\begin{tabular}{|l|c|}
\hline
\textbf{Target Node}  & \textbf{CRS Applied Count} \\
\hline\hline
Node 2 (Math Analyst)  & 228 \\
\hline
Node 3 (Programmer)    & 212 \\
\hline
Node 0 (FinalRefer)    & 184 \\
\hline\hline
\textbf{Total}         & \textbf{624} \\
\hline
\end{tabular}
\end{table}
